\chapter{Bases de Python pour les professionnels de la sant\'e}
\label{ch:bases-python}

\begin{quote}
\emph{<<\,L'objectif n'est pas d'apprendre Python. L'objectif est de r\'epondre \`a des questions de sant\'e avec des donn\'ees. Python n'est que l'outil.\,>>}
\end{quote}

% ============================================================
\section{Pourquoi Python pour les donn\'ees de sant\'e ?}

Trois raisons pour lesquelles les professionnels de la sant\'e devraient apprendre Python plut\^ot que de se fier uniquement \`a Excel ou SPSS :

\begin{enumerate}
  \item \textbf{Reproductibilit\'e.} Un script Python documente chaque \'etape de votre analyse. Quand un examinateur demande <<\,comment avez-vous trait\'e les valeurs de laboratoire manquantes ?\,>>, vous montrez le code --- pas une description de clics dans des menus.
  \item \textbf{\'Echelle.} Excel a du mal avec 100\,000 dossiers de patients. Python en g\`ere des millions sans ralentir.
  \item \textbf{\'Ecosyst\`eme.} Des biblioth\`eques comme \texttt{lifelines} (analyse de survie), \texttt{geopandas} (cartographie des maladies) et \texttt{scikit-learn} (mod\`eles pr\'edictifs) offrent des outils que SPSS ne propose pas.
\end{enumerate}

\begin{clinicalnote}
Vous n'aurez pas besoin de m\'emoriser la syntaxe. Les professionnels de la sant\'e utilisent Python comme les cliniciens utilisent un st\'ethoscope --- vous apprenez les parties dont vous avez besoin et vous cherchez le reste. Chaque bloc de code de ce cours peut \^etre copi\'e dans un notebook Jupyter et ex\'ecut\'e imm\'ediatement.
\end{clinicalnote}

% ============================================================
\section{Configurer votre environnement}

\subsection{Google Colab (recommand\'e)}

Aucune installation. Rendez-vous sur \url{https://colab.research.google.com}, connectez-vous avec un compte Google et cr\'eez un nouveau notebook. Toutes les biblioth\`eques que nous utilisons sont pr\'einstall\'ees.

\subsection{Installation locale avec Anaconda}

Si vous pr\'ef\'erez travailler hors ligne :
\begin{enumerate}
  \item T\'el\'echargez Anaconda : \url{https://www.anaconda.com/download}
  \item Installez avec les param\`etres par d\'efaut.
  \item Ouvrez \textbf{Jupyter Notebook} depuis Anaconda Navigator.
\end{enumerate}

% ============================================================
\section{Votre premier notebook}

Un notebook Jupyter est un document qui m\'elange texte, code et r\'esultats. Chaque \emph{cellule} contient soit :
\begin{itemize}
  \item \textbf{Du code} --- des instructions Python qui produisent un r\'esultat.
  \item \textbf{Du Markdown} --- du texte format\'e (explications, titres, notes).
\end{itemize}

Appuyez sur \texttt{Shift + Entr\'ee} pour ex\'ecuter une cellule.

\begin{minted}{python}
# Votre premiere cellule Python
print("Bonjour, donnees de sante !")
\end{minted}

% ============================================================
\section{Variables et types}

Une \emph{variable} stocke une valeur. En donn\'ees de sant\'e, vous travaillez constamment avec quatre types :

\begin{minted}{python}
patient_age = 45              # int (entier)
temperature = 37.2            # float (nombre decimal)
diagnosis = "Type 2 Diabetes" # str (chaine de caracteres / texte)
is_hospitalized = True        # bool (True ou False)
\end{minted}

\begin{pythontip}
Les noms de variables doivent \^etre descriptifs. Utilisez \texttt{patient\_age}, pas \texttt{x}. Votre futur vous (et vos collaborateurs) vous remercieront.
\end{pythontip}

% ============================================================
\section{Collections : listes et dictionnaires}

\subsection{Listes}

Une liste contient plusieurs valeurs dans un ordre donn\'e :
\begin{minted}{python}
blood_pressures = [120, 135, 128, 142, 118]
print(f"Nombre de mesures : {len(blood_pressures)}")
print(f"Premiere mesure : {blood_pressures[0]}")  # l'indexation commence \`a 0
print(f"Derniere mesure : {blood_pressures[-1]}")
\end{minted}

\subsection{Dictionnaires}

Un dictionnaire associe des cl\'es \`a des valeurs --- comme un dossier patient :
\begin{minted}{python}
patient = {
    "id": "PAT-0042",
    "age": 58,
    "sex": "F",
    "diagnosis": "Hypertension",
    "systolic_bp": 148,
    "medications": ["Amlodipine", "Hydrochlorothiazide"]
}
print(patient["diagnosis"])  # Hypertension
\end{minted}

% ============================================================
\section{Structures de contr\^ole}

\subsection{Conditions}

\begin{minted}{python}
systolic = 148

if systolic >= 140:
    category = "Hypertension de stade 2"
elif systolic >= 130:
    category = "Hypertension de stade 1"
elif systolic >= 120:
    category = "Elevee"
else:
    category = "Normale"

print(f"Categorie de PA : {category}")
\end{minted}

\subsection{Boucles}

\begin{minted}{python}
patients = [
    {"id": "P001", "age": 34, "hba1c": 5.4},
    {"id": "P002", "age": 67, "hba1c": 7.8},
    {"id": "P003", "age": 52, "hba1c": 6.1},
]

for p in patients:
    if p["hba1c"] >= 6.5:
        print(f"{p['id']}: Diabetique (HbA1c = {p['hba1c']})")
    elif p["hba1c"] >= 5.7:
        print(f"{p['id']}: Pre-diabetique (HbA1c = {p['hba1c']})")
    else:
        print(f"{p['id']}: Normal (HbA1c = {p['hba1c']})")
\end{minted}

% ============================================================
\section{Fonctions}

Une fonction encapsule un calcul r\'eutilisable :

\begin{minted}{python}
def bmi(weight_kg, height_m):
    """Calculer l'indice de masse corporelle."""
    return weight_kg / (height_m ** 2)

def bmi_category(bmi_value):
    """Classifier l'IMC selon les categories de l'OMS."""
    if bmi_value < 18.5:
        return "Insuffisance ponderale"
    elif bmi_value < 25.0:
        return "Poids normal"
    elif bmi_value < 30.0:
        return "Surpoids"
    else:
        return "Obesite"

value = bmi(82, 1.75)
print(f"IMC : {value:.1f} ({bmi_category(value)})")
\end{minted}

\begin{exercise}
\begin{enumerate}
  \item \'Ecrivez une fonction \texttt{classify\_bp(systolic, diastolic)} qui retourne une cat\'egorie de pression art\'erielle selon les recommandations AHA/ACC.
  \item Cr\'eez une liste de 5 dictionnaires patients avec les champs \texttt{weight\_kg} et \texttt{height\_m}. Parcourez-les, calculez l'IMC et affichez la cat\'egorie pour chacun.
  \item \'Ecrivez une fonction \texttt{egfr(creatinine, age, sex)} qui estime le d\'ebit de filtration glom\'erulaire selon l'\'equation CKD-EPI.
\end{enumerate}
\end{exercise}

% ============================================================
\section{Importer des biblioth\`eques}

La puissance de Python vient de ses biblioth\`eques. Voici celles que nous utiliserons tout au long de ce cours :

\begin{minted}{python}
import pandas as pd          # manipulation de donnees
import numpy as np           # calcul numerique
import matplotlib.pyplot as plt  # graphiques
import seaborn as sns        # visualisation statistique
from scipy import stats      # tests statistiques
\end{minted}

\begin{warningbox}
Si vous obtenez \texttt{ModuleNotFoundError}, installez la biblioth\`eque manquante :
\begin{minted}{bash}
pip install pandas matplotlib seaborn scipy
\end{minted}
Dans Google Colab, pr\'efixez avec \texttt{!} : \texttt{!pip install lifelines}
\end{warningbox}

% ============================================================
\section{R\'esum\'e du chapitre}

\begin{itemize}
  \item Python est gratuit, reproductible et dispose de biblioth\`eques sp\'ecifiques \`a la sant\'e.
  \item Les notebooks Jupyter m\'elangent code, texte et r\'esultats dans un seul document.
  \item Les variables stockent des valeurs ; les listes et dictionnaires les organisent.
  \item Les fonctions rendent votre analyse r\'eutilisable et lisible.
  \item \texttt{pandas}, \texttt{matplotlib}, \texttt{seaborn} et \texttt{scipy} sont les biblioth\`eques fondamentales pour l'analyse de donn\'ees de sant\'e.
\end{itemize}
